\chapter*{Conclusion}
\addcontentsline{toc}{chapter}{Conclusion}

The aim of this thesis was to understand, study, analyze, and implement
classical knapsack cryptosystems. The thesis did not focus only on the original
Merkle--Hellman variant, but also analyzed and compared selected related
variants by looking at their design, efficiency, and known weaknesses. The main
motivation was to show why the difficulty of the underlying mathematical
problem is not enough to assume that a cryptosystem is secure, especially when
the trapdoor is based on an easy structured version of that problem.

The theoretical part firstly introduced cryptographic and complexity-theory
background needed for the topic. Special focus was given to the distinction
between general subset-sum hardness and the structured instances used in
cryptosystems. It then described the Merkle--Hellman construction, including
the role of the superincreasing private sequence, how modular multiplication
creates the public key, and how encryption and decryption work. It then
described subset-sum as the underlying mathematical problem for the
Merkle--Hellman cryptosystem. After defining the problem, it explained that the
decision version of subset-sum is NP-complete and why this fact was important
for early knapsack cryptography. The role of the superincreasing private
sequence and modular multiplication was explained before moving to the
individual phases of the system: key generation, encryption, and decryption.
The later sections discussed why the basic construction does not hide its
trapdoor structure well enough and is therefore vulnerable. The main attacks
described were Shamir's structural attack and low-density lattice-based
attacks. The thesis also described selected variants, including permuted and
iterated Merkle--Hellman constructions and the Chor--Rivest knapsack-type
cryptosystem, together with their known cryptanalysis.

The practical part of the thesis focused on implementing the classical
Merkle--Hellman cryptosystem and selected variants in modular C
program. The implementation was not intended to be a secure cryptographic
library, but rather an experimental tool with separate modes for demonstration,
benchmarking, and simple solver experiments.

The program implements complete procces with key-generation, encryption and
decryption. The benchmark mode measured key generation, encryption, and
decryption times and recorded values such as density, modulus margin, and the
bit length of the private-weight sum. The solver experiments tested direct
subset-sum methods such as brute force and meet-in-the-middle search on small
generated instances. 

During implementation, it also made several design
choices that are often left open in high-level descriptions of the
cryptosystem. For example how the superincreasing sequence is generated, how
much margin is added before choosing the modulus, and how the multiplier is
sampled.

The experiments showed that key generation dominates the runtime of the
implemented schemes for the tested block sizes. The permuted variant behaves
similarly to the classical variant, while additional iterated layers noticeably
increase the cost. The iterated-layer sweep also showed that adding modular
layers lowers the measured density of the public knapsack, especially for small
block sizes. This agrees with the theoretical warning that iteration does not
remove the fundamental weakness and may even move the construction toward
size-related properties associated with low-density attacks.

The parameter sweeps showed that the details of the private sequence generator
matter. Increasing the random increment bound makes the private sum and public
weights slightly larger and therefore lowers the measured density, most visibly
for smaller block sizes. The modulus margin mainly demonstrates the necessary
condition for decryption: the modulus must be larger than every possible private
subset sum. These results do not measure security directly, but they connect the
implementation to the density discussion from the theoretical part.

The brute-force and meet-in-the-middle experiments illustrated the growth of
direct subset-sum solving methods on small instances. Meet-in-the-middle extends
the feasible range compared with brute force, but it still requires exponential
memory. These solvers are therefore useful as educational baselines, while the
real cryptanalytic weakness of Merkle--Hellman comes from structural attacks
such as Shamir's attack and from lattice-based methods.

Overall, the work supports the main lesson of classical knapsack cryptography:
a hard mathematical problem does not automatically produce a secure public-key
cryptosystem. The trapdoor construction, the distribution of generated public
keys, and the available cryptanalytic methods are decisive. Merkle--Hellman and
its variants are no longer suitable for practical encryption, but they remain a
valuable historical example because they show clearly how implementation
choices, parameter choices, and hidden structure interact in public-key
cryptography.

Possible extensions of this work would include implementing Shamir's structural
attack or an LLL-based low-density solver, adding more historical variants, or
measuring key and ciphertext sizes in more detail. These additions would make
the experimental comparison closer to the cryptanalytic history, but they would
not change the main conclusion: classical trapdoor knapsack cryptosystems are
best treated as educational systems rather than secure cryptographic tools.
